\documentclass[11pt]{amsart}
\usepackage{calc,amssymb,amsmath,ulem,stmaryrd,fullpage}
\usepackage{enumerate}
\usepackage{alltt}
\RequirePackage[dvipsnames,usenames]{color}
\let\mffont=\sf
\normalem
%\input{kmacros3.sty}
\input{xy}
\xyoption{all}
\usepackage{hyperref}
\usepackage{graphicx}
\usepackage[all,cmtip]{xy}
\usepackage{verbatim}
\usepackage[left=0.95in,top=0.95in,right=0.95in,bottom=0.95in]{geometry}
\newcommand{\tauCohomology}{T}
\newcommand{\FCohomology}{S}


\theoremstyle{theorem}
%\newtheorem*{mainthm*}{Main Theorem}

\newcommand{\note}[1]{\marginpar{\sffamily\tiny #1}}

\def\todo#1{\textcolor{Mahogany}%
{\sffamily\footnotesize\newline{\color{Mahogany}\fbox{\parbox{\textwidth-15pt}{#1}}}\newline}}

\renewcommand{\O}{\mathcal O}
\newcommand{\ie}{{\itshape i.e.} }
\newcommand{\roundup}[1]{\lceil #1 \rceil}
\newcommand{\rounddown}[1]{\lfloor #1 \rfloor}
\renewcommand{\to}[1][]{\xrightarrow{\ #1\ }}
\newcommand{\onto}[1][]{\protect{\xrightarrow{\ #1\ }\hspace{-0.8em}\rightarrow}}
\newcommand{\into}[1][]{\lhook \joinrel \xrightarrow{\ #1\ }}
%\newcommand{DBDual}[1]{{\underline{\omega}_{#1}}}
\newcommand{\DBDual}[1]{{{\uuline \omega}{}^{\mydot}_{#1}}}
\newcommand{\Tt}{{\mathfrak{T}}}
%\newcommand{\bn}{{\mathfrak{n}} }
\newcommand{\sol}[1]{ \vskip 6pt{\color{blue} {\bf Solution:} #1} \vskip 6pt}

\newcommand{\bR}{{\mathbb{R}}}
\newcommand{\va}{{\vec{a}}}
\newcommand{\vb}{{\vec{b}}}
\newcommand{\vzero}{{\vec{0}}}
\newcommand{\vi}{{\vec{i}}}
\newcommand{\vj}{{\vec{j}}}
\newcommand{\vk}{{\vec{k}}}
\newcommand{\vh}{{\vec{h}}}
\newcommand{\vr}{{\vec{r}}}
\newcommand{\vu}{{\vec{u}}}
\newcommand{\vv}{{\vec{v}}}
\newcommand{\vw}{{\vec{w}}}
\newcommand{\vx}{{\vec{x}}}
\newcommand{\vy}{{\vec{y}}}
\newcommand{\vz}{{\vec{z}}}
\newcommand{\vT}{{\vec{T}}}
\newcommand{\vN}{{\vec{N}}}
\newcommand{\vF}{{\vec{F}}}
\newcommand{\vG}{{\vec{G}}}
\newcommand{\bF}{\mathbb{F}}
\newcommand{\bQ}{\mathbb{Q}}
\newcommand{\bZ}{\mathbb{Z}}
\newcommand{\Inn}{\mathrm{Inn}}
\newcommand{\Aut}{\mathrm{Aut}}
\newcommand{\Ann}{\mathrm{Ann}}
\newcommand{\tor}{\mathrm{tor}}


\begin{document}
\title{Worksheet \#5 -- Math 6310\\Fall 2019}
\author{Due: Wednesday, November 6th}
\maketitle

\noindent
\vskip 9pt
\noindent
{\bf 1.}  $R = \bZ$.  Consider the matrix
\[M = 
\left[ \begin{array}{ccc}  
1 & 2 & 3 \\
4 & 5 & 6 \\
7 & -8 & 9
\end{array}  \right]
\]
Write the cokernel of $R^3 \xrightarrow{M \cdot } R^3$ as a direct sum of cyclic modules.
\newpage
\noindent
{\bf 2.}  Consider $R = \bQ[x]$ and the matrix:
\[ M = 
\left[ \begin{array}{cc}
x & 0  \\
x & x^2  \\
1 & 1 
\end{array}  \right]
\]
Write the cokernel of $R^2 \xrightarrow{M \cdot} R^3$ as a direct sum of cyclic modules.
\newpage
\noindent
{\bf 3.}  Consider $R = \bZ[i]$ with matrix 
\[ M = 
\left[ \begin{array}{ccc}
2 & i & -1 \\
0 & 2i & 1+i \\
i & 1 & 2
\end{array}
\right]
\]
Write the cokernel of $R^3 \xrightarrow{M \cdot} R^3$ as a direct sum of cyclic modules.
\newpage
\noindent
{\bf 4.}  Suppose that $R = \mathbb{C}[x]$, and that $M$ is the $R$-module generated by three elements $a, b$, and $c$, modulo the three relations $a-xc$, \ $x a - xb + xc$, and $xb - (x^2-1)c$.  Write $M$ as a direct sum of cyclic modules.
\newpage
\noindent
{\bf 5.}  Let $A =\bF_3[x]$, i.e., a polynomial ring in one variable over the field with 3 elements.  Suppose $M$ and $N$ are finitely generated $A$-modules such that
\[
M \oplus {A \over {x^3 + 1}}\ \cong\ N \oplus {A \over {x + 1}} \oplus {A \over {x + 1}} \oplus {A \over {x+1}}.
\]
Are $M$ and $N$ isomorphic?  Justify your answer.
\newpage
\noindent
{\bf 6.}  Let $R = \bQ[x]$ and consider the submodule $M$ of $R^2$ generated by the elements $(x^2-1,\ x-1)$ and $(x^2+x,\ x)$.  Write $M$ as a direct sum of cyclic modules.
\newpage
\noindent
{\bf 7.}  Let $M$ be the cokernel of the mapping from $\mathbb{Z}^2$ to $\mathbb{Z}^3$ given by the matrix
$$
\begin{bmatrix}
2 & 8\\ 
4 & 10\\
6 & 12
\end{bmatrix}
$$
How many $\mathbb{Z}$-module homomorphisms are there from $M$ to $\mathbb{Z}/3\mathbb{Z}$?

\end{document}

